\subsection{测试函数}

在本章中，U 表示 \(\mathbf{R}^{n}\) 的开子集。对 \(p \in [1, +\infty]\) ，我们将写 \(\mathcal{L}^{p}(U)\) 与 \(L^{p}(U)\) ，而不写 \(\mathcal{L}^{p}(U, \mu)\) 与 \(L^{p}(U, \mu)\) 。我们把属于 \(\mathcal{L}^{p}(U)\) 的连续函数与其在 \(L^{p}(U)\) 中的像等同起来。

给 U 上的复值无穷次可微函数空间 \(\mathcal{C}^{\infty}(\mathrm{U})\) 赋以由半范数族 \(p_{\alpha,K}\) ——对 \(\alpha \in \mathbf{N}^{n}\) 和 \(K \subset U\) 定义为

\[
p _ {\alpha , \mathrm{K}} (\varphi) = \sup _ {x \in \mathrm{K}} | \partial^ {\alpha} \varphi (x) |.
\]

该空间是完备的（参见 EVT, III, p. 9，例 b））。

对 U 的每个紧子集 K，我们用 \(\mathcal{C}_{\mathrm{K}}^{\infty}(\mathrm{U})\) 表示 \(\mathcal{C}^{\infty}(\mathrm{U})\) 中由支集含于 K 的函数构成的子空间。空间 \(\mathcal{C}_{\mathrm{K}}^{\infty}(\mathrm{U})\) 赋以由 \(\mathcal{C}^{\infty}(\mathrm{U})\) 的拓扑诱导的拓扑。它是一个弗雷歇空间，定义其拓扑的可数半范数族即族 \((p_{\alpha,\mathrm{K}})_{\alpha\in\mathbf{N}^{n}}\) 。特别地，它是有界型空间（EVT, III, p. 12，命题 2）。

我们用 \(\mathcal{D}(\mathrm{U})\) 表示向量空间 \(\mathcal{K}(\mathrm{U})\cap \mathcal{C}^{\infty}(\mathrm{U})\) ——由 U 中具有紧支集的 \(\mathrm{C}^{\infty}\) 类函数所构成。我们称 \(\mathcal{D}(\mathrm{U})\) 为 U 上的测试函数空间。空间 \(\mathcal{D}(\mathrm{U})\) 是诸空间 \(\mathcal{C}_{\mathrm{K}}^{\infty}(\mathrm{U})\) 的归纳极限，并赋以相应的局部凸归纳极限拓扑（EVT, II, p. 31）。

该空间在 EVT, III, p. 9 中记作 \(\mathcal{C}_{\circ}^{\infty}(U)\) 。

设 \((\mathrm{K}_{m})\) 是 U 的紧子集的递增序列，其内部构成 U 的覆盖。那么空间 \(\mathcal{D}(\mathrm{U})\) 是诸空间 \(\mathcal{C}_{\mathrm{K}_{m}}^{\infty}(\mathrm{U})\) 的严格归纳极限（EVT, III, p. 9）。因此它是完备空间（EVT, II, p. 35，命题 9）。\(\mathcal{D}(\mathrm{U})\) 的每个有界子集都含于某个子空间 \(\mathcal{C}_{\mathrm{K}_{m}}^{\infty}(\mathrm{U})\) 之中（EVT, III, p. 5，命题 6）。这就是说，\(\mathrm{B}\subset\mathcal{D}(\mathrm{U})\) 有界当且仅当存在 U 的紧子集 K 以及族 \((\mathrm{M}_{\alpha})_{\alpha\in\mathbf{N}^{n}}\) （\(R_{+}\) 中的），使得 B 含于诸函数 \(\varphi\in\mathcal{C}_{\mathrm{K}}^{\infty}(\mathrm{U})\) ——满足

\[
p _ {\alpha , \mathrm{K}} (\varphi) \leqslant \mathrm{M} _ {\alpha}
\]

——所成之集（对所有 \(\alpha\in\mathbf{N}^{n}\) ）。

空间 \(\mathcal{D}(\mathrm{U})\) 是有界型空间（EVT, III, p. 12，例 3）和蒙泰尔空间（EVT, IV, p. 18，例 4）。

设 V 是含于 U 的 \(\mathbf{R}^{n}\) 的开子集。若 \(K \subset V\) 是紧的，则函数到 V 上的限制定义了从 \(\mathcal{C}_{\mathrm{K}}^{\infty}(\mathrm{U})\) 到 \(\mathcal{C}_{\mathrm{K}}^{\infty}(\mathrm{V})\) 上的连续满线性映射。定义在 V 上的函数用零延拓，诱导一个从 \(\mathcal{D}(\mathrm{V})\) 到 \(\mathcal{D}(\mathrm{U})\) 中的单连续线性映射，称为典范映射。

注. --- 我们类似地定义 U 上的实值测试函数空间 \(\mathcal{D}_{\mathbf{R}}(\mathrm{U})\) 。使得 \(z \otimes \varphi \mapsto z\varphi\) （对所有 \(z \in \mathbf{C}\) 及每个 \(\varphi \in \mathcal{D}_{\mathbf{R}}(\mathrm{U})\) ）的线性映射是 \(\mathbf{C} \otimes \mathcal{D}_{\mathbf{R}}(\mathrm{U})\) 到 \(\mathcal{D}(\mathrm{U})\) 上的同构。

引理 2. --- 设 \(\mathcal{U}\) 是 U 的开覆盖。存在比 \(\mathcal{U}\) 更细的、由相对紧集合组成的 U 的局部有限开覆盖。

由于 U 是局部紧的，存在 \(\mathcal{V}\) ——U 的覆盖，比 \(\mathcal{U}\) 更细，由在 U 中相对紧的开集组成。由于 U 是仿紧的（TG, IX, p. 51，定理 4），存在局部有限开覆盖 \(\mathcal{W}\) ，比 \(\mathcal{V}\) 更细。于是 \(\mathcal{W}\) 比 \(\mathcal{U}\) 更细，且由相对紧的开集组成。

引理 3. --- 设 \(f \in \mathcal{K}(\mathbf{R}^n)\) ，V 是 f 的支集 K 的开邻域。存在整数 \(m_0 \geqslant 1\) ，使得对每个 \(x \in K\) ，集合 V 包含以 x 为中心、半径为 \(m_0^{-1}\) 的球。对每个整数 \(m > m_0\) ，设 \(V_m\) 是以 0 为中心、半径为 \(m^{-1}\) 的闭球（在 \(\mathbf{R}^{n}\) 中），\(\varphi_{m}\) 是支集含于 \(V_{m}\) 且积分为 1 的正测试函数。令 \(f_{m} = \varphi_{m} * f\) 。

\begin{enumerate}
\def\labelenumi{\alph{enumi})}
\item
  有 \(f_{m}\in \mathcal{D}(\mathbf{R}^{n})\) ，且 \(f_{m}\) 的支集含于 V；
\item
  设 \(p \in [1, +\infty]\) 。序列 \((f_m)_{m > m_0}\) 收敛于 \(f\) （在 \(L^p(\mathbf{R}^n)\) 中）。
\end{enumerate}

依据 TG, IX, p. 14，注，我们有 \(d(\mathrm{K}, \mathbf{R}^{n}-\mathrm{V}) > 0\) 。设 \(m_{0} \geqslant 1\) 是满足 \(m_{0}^{-1} < d(\mathrm{K}, \mathbf{R}^{n}-\mathrm{V})\) 的整数；它满足所需的条件。对所有 \(m > m_{0}\) ，函数 \(f_{m}\) 连续（INT, VIII, p. 166, § 4, n° 5，命题 11），且其支集含于 \(\mathrm{K} + \mathrm{V}_{m}\) （INT, VIII, p. 126, § 1, n° 4，命题 5），从而含于 V。依据 IV, p. 198 的推论 2，我们有 \(f_{m} \in \mathcal{D}(\mathbf{R}^{n})\) 。若 \(p \neq +\infty\) ，则序列 ( \(f_{m}\) ) 收敛于 \(f\) （在 \(\mathrm{L}^{p}(\mathbf{R}^{n})\) 中；INT, VIII, p. 172, § 4, no. 7，命题 20）。\(^{(1)}\)

设 \(p = +\infty\) 。设 \(\varepsilon > 0\) 。函数 \(f\) 在 \(\mathbf{R}^n\) 上一致连续；因此存在整数 \(m_1 > m_0\) ，使得对每个 \(m \geqslant m_1\) 以及所有 \(x \in \mathbf{R}^n\) 和 \(y \in \mathrm{V}_m\) ，有 \(|f(x) - f(x - y)| < \varepsilon\) 。

对每个 \(m \geqslant m_{1}\) ，函数 \(\varphi_{m}\) 在 \(\mathrm{V}_{m}\) 外为零、是正的且积分为 1。因此我们推出不等式

\[
| f (x) - f _ {m} (x) | = \left| \int_ {\mathbf {R} ^ {n}} (f (x) - f (x - y)) \varphi_ {m} (y) d \mu (y) \right| \leqslant \varepsilon
\]

（对所有 \(x \in \mathbf{R}^{n}\) ），由此得结论。

命题 4. --- a) \(\mathcal{D}(\mathrm{U})\) 到 \(\mathcal{K}(\mathrm{U})\) 中的典范单射是连续的，且 \(\mathcal{D}(\mathrm{U})\) 在 \(\mathcal{K}(\mathrm{U})\) 中稠密；

\begin{enumerate}
\def\labelenumi{\alph{enumi})}
\setcounter{enumi}{1}
\tightlist
\item
  对每个 \(p \in [1, +\infty[\) ，\(\mathcal{D}(\mathrm{U})\) 到 \(\mathcal{L}^p(\mathrm{U})\) 中的典范单射是连续的，且 \(\mathcal{D}(\mathrm{U})\) 在 \(\mathrm{L}^p(\mathrm{U})\) 中稠密。
\end{enumerate}

\(\mathcal{K}(\mathrm{U})\) 上的每个连续半范数在 \(\mathcal{D}(\mathrm{U})\) 上连续，因而 \(\mathcal{D}(\mathrm{U})\) 到 \(\mathcal{K}(\mathrm{U})\) 中的典范单射是连续的。

设 \(f \in \mathcal{K}(\mathrm{U})\) 。存在序列 \((f_{m})_{m \in \mathbf{N}}\) （在 \(\mathcal{D}(\mathrm{U})\) 中），它在 U 上一致收敛于 f，且诸 \(f_{m}\) 的支集都含于 f 的支集的一个固定的相对紧邻域（引理 3）。因此序列 \((f_{m})\) 收敛于 f （在 \(\mathcal{K}(\mathrm{U})\) 中）。这就完成了 a) 的证明。

设 \(p \in [1, +\infty[\) 。设 \(f \in \mathcal{D}(U)\) ，K 是 f 的支集。我们于是有不等式 \(\mathrm{N}_{p}(f) \leqslant \mu(\mathrm{K})^{1/p} \sup_{x \in \mathrm{K}} |f(x)|\) ，因而 \(\mathcal{D}(U)\) 到 \(\mathcal{L}^{p}(U)\) 中的典范单射是连续的。断言 b) 的后半部分由 a) 得出。

两个测试函数的乘积仍是测试函数，故莱布尼茨公式（FVR, I, p. 28，命题 2）表明空间 \(\mathcal{D}(\mathrm{U})\) 是一个拓扑代数。

更一般地，设 \(f \in \mathcal{C}^{\infty}(\mathrm{U})\) 。那么线性映射 \(\varphi \mapsto f\varphi\) ——从 \(\mathcal{D}(\mathrm{U})\) 到 \(\mathcal{D}(\mathrm{U})\) 中——是连续的。事实上，对每个紧子集 \(\mathbf{K}\) （\(\mathrm{U}\) 的）及每个 \(\alpha \in \mathbf{N}^{n}\) ，我们有

\[
p _ {\alpha , \mathrm{K}} (f \varphi) \leqslant \sum_ {0 \leqslant \beta \leqslant \alpha} \binom {\alpha} {\beta} p _ {\beta , \mathrm{K}} (f) p _ {\alpha - \beta , \mathrm{K}} (\varphi)
\]

（参见同前）。

